\chapter{Label classes in the simple case}
\label{ch:chain-simple}

We begin with offspring laws supported on $\set{0,1,2}$. In the chain regime, the branch
points form a binary tree and the intervening chain lengths become independent labels. We
quantise those labels so that the matching theorem of \cref{ch:iid} applies. In the bushy
regime, the Harris decomposition replaces individual chain lengths by finite decorated shapes,
to which the same strategy applies. These results are formalised in \texttt{ChainClasses} over
the Galton--Watson foundation \texttt{BranchingProcess}.

\section{The chain encoding}

Assume first $\theta_0=0$, so that $\theta_1+\theta_2=1$. Every vertex of $\cT$ then lies on a
unique maximal chain of degree-two vertices ending at a branch point, and the branch points
form a copy of the binary tree $\bbB$. Write $\iota(w)$ for the vertex of $\cT$ carrying the
branch point $w\in\bbB$, and $\lambda(w)$ for the length of the chain issuing from it.

\begin{lemma}[The tree is the disjoint union of its chains]
\label{thm:chains}
\lean{ChainClasses.tree_eq_chains}
\leanok
\uses{def:gw}
Every $v\in\cT$ has a unique representation $v=\iota(w)1^{l}$ with $w\in\bbB$ and
$0\leq l\leq\lambda(w)-1$. If $w$ is not a prefix of $w'$, then the ray
$\set{\iota(w)1^{j}\colon j\geq0}$ is disjoint from the chain of $w'$.
\end{lemma}

\begin{proof}
\leanok
The chain recursion appends letters without changing any previously constructed address. Thus
prefixes in $\bbB$ map to prefixes in $\cT$. Given $v\in\cT$, follow the path from the root to
$v$ and choose its last branch point, say $\iota(w)$. The remaining part of the path consists
only of first children, so $v=\iota(w)1^l$ for a unique offset
$0\leq l<\lambda(w)$. This proves existence and uniqueness. If $w$ is not a prefix of $w'$,
the paths to the corresponding branch points separate below their last common ancestor, so the
stated ray cannot meet the chain of $w'$.
\end{proof}

\begin{lemma}[The chain lengths are i.i.d.\ geometric]
\label{thm:geometric}
\lean{ChainClasses.label_iIndepFun}
\leanok
The random variables $\lambda(w)$, $w\in\bbB$, are independent and geometrically distributed,
with $\bbP(\lambda(w)=n)=\theta_1^{n-1}\theta_2$ for $n\geq1$.
\end{lemma}

\begin{proof}
\leanok
\uses{thm:chains}
Fix a finite prefix-closed set $F\subseteq\bbB$ and prescribe lengths $n_w\geq1$ for
$w\in F$. Explore the associated chains in prefix order. Conditional on everything already
revealed, the chain at $w$ consists of $n_w-1$ unary vertices followed by one binary vertex.
Its conditional probability is therefore $\theta_1^{n_w-1}\theta_2$, independently of the
earlier chains. Multiplying these conditional probabilities gives
\[
 \bbP\bigl(\lambda(w)=n_w\text{ for all }w\in F\bigr)
 =\prod_{w\in F}\theta_1^{n_w-1}\theta_2.
\]
The finite-dimensional distributions consequently have the asserted independent geometric
form.
\end{proof}

\begin{lemma}[Relabelling is an isometry]
\label{thm:isometry}
\lean{ChainClasses.isometry_of_relabel}
\leanok
Let $\pi\in\Aut(\bbB)$ and let $\lambda$ be a labelling with associated tree $\cT$. The tree
$\cT'$ associated with $\lambda'(w)=\lambda(\pi^{-1}w)$ is isometric to $\cT$.
\end{lemma}

\begin{proof}
\leanok
\uses{thm:chains}
By \cref{thm:chains}, every vertex of $\cT$ has a unique form $\iota(w)1^l$. Send it to the
vertex at the same offset $l$ in the chain indexed by $\pi(w)$ in $\cT'$. The definition of
$\lambda'$ ensures that the two chains have the same length, so this map is a bijection. An
automorphism of $\bbB$ preserves prefixes and last common ancestors. The map therefore preserves
both distances within a chain and the sums of chain lengths between distinct chains, hence is an
isometry.
\end{proof}

\begin{lemma}[Transfer]
\label{thm:transfer}
\lean{ChainClasses.transfer}
\leanok
\uses{thm:chains, def:qi}
Let $\lambda,\lambda'$ be labellings of $\bbB$ with values in $\bbN$ and let $D\geq1$ satisfy
$D^{-1}\lambda'(w)\leq\lambda(w)\leq D\,\lambda'(w)$ for all $w\in\bbB$. Then there is a
$(D+3)$-quasi-isometry $\psi\colon\cT\to\cT'$ between the associated trees.
\end{lemma}

\begin{proof}
\leanok
\uses{thm:chains}
For each $w$, map the chain of length $\lambda(w)$ monotonically to the chain of length
$\lambda'(w)$ by proportional division, fixing the branch point. The hypothesis bounds the
distortion of each individual chain by $D$. If two vertices lie in different chains, their
geodesic is a concatenation of two partial chains and the complete chains indexed by the words
between them and their last common ancestor. Applying the same comparison to every term gives
the required upper and lower distance bounds. Consecutive image points on a target chain are at
distance at most $D+2$, so every target vertex lies within the required coarse-density bound.
Combining the distortion and density estimates gives a $(D+3)$-quasi-isometry.
\end{proof}

\section{Quantisation into level classes}

Write $\ell_D(n):=\floor{\log_D n}$ for the level of $n\geq1$ at scale $D$.

\begin{lemma}[Neighbouring levels are comparable]
\label{thm:level}
\lean{ChainClasses.level_close_comparable}
\leanok
Let $D\geq2$ and $m,m'\in\bbN$. If $\abs{\ell_D(m)-\ell_D(m')}\leq1$ then $D^{-2}m'<m<D^2m'$.
\end{lemma}

\begin{proof}
\leanok
With $j=\ell_D(m)$ and $j'=\ell_D(m')$ we have $D^j\leq m<D^{j+1}$ and likewise for
$m'$. From $j\leq j'+1$ it follows that $m<D^{j'+2}\leq D^2m'$, and the lower bound is the
same argument with the roles exchanged.
\end{proof}

\begin{lemma}[The quantised law]
\label{thm:quantised-law}
\lean{ChainClasses.qF_le}
\leanok
Let $p^{(D)}$ be the image under $\ell_D$ of the geometric law $\theta_1^{n-1}\theta_2$,
$n\geq1$. Then $p^{(D)}_0=1-\theta_1^{D-1}$ and $p^{(D)}_j\leq\theta_1^{D^j-1}$ for $j\geq1$.
\end{lemma}

\begin{proof}
\leanok
\uses{thm:geometric}
The geometric law has tail $\theta_1^{m-1}$ at $m$, and $p^{(D)}_j$ is the difference of
the tails at $D^j$ and $D^{j+1}$.
\end{proof}

\begin{lemma}[Chain classes]
\label{thm:chain-classes}
\lean{ChainClasses.quantised_label_iIndepFun}
\leanok
Let $\theta_1+\theta_2=1$ with $0<\theta_1<1$, let $D\geq2$ be an integer, and put
$\cC_k=\set{n\in\bbN\colon D^k\leq n<D^{k+1}}$. The classes $\cC_k$ partition $\bbN$, any two
$n,m\in\cC_k$ satisfy $D^{-1}<n/m<D$, and the quantised labels $\ell_D(\lambda(w))$,
$w\in\bbB$, are independent with common law
\[
 p^{(D)}_k=\bbP\bigl(\lambda(w)\in\cC_k\bigr)=\theta_1^{D^k-1}-\theta_1^{D^{k+1}-1}.
\]
\end{lemma}

\begin{proof}
\leanok
\uses{thm:geometric, thm:quantised-law}
The classes are the fibres of $\ell_D$, hence a partition of $\bbN$, and dividing the
defining inequalities by $D^k$ gives the two ratio bounds. Independence is inherited from the
labels, and the class masses are the tail differences of the geometric law.
\end{proof}

\begin{lemma}[Chain coupling across two laws]
\label{thm:chain-coupling}
\lean{ChainClasses.coupling_law}
\leanok
Let $\theta_1,\theta_1'\in(0,1)$, put $\gamma=\log\theta_1/\log\theta_1'>0$, and let $D\geq2$
be an integer with $\gamma(D-1)\geq1$. Let $p^{(D)}$ be the law of \cref{thm:chain-classes} for
$\theta_1$ and $D$, let $\lambda'$ be geometric with parameter $\theta_1'$, and let $U$ be
uniform on $[0,1)$ and independent of $\lambda'$. Then there are cut points
$1=D_0'<D_1'<D_2'<\cdots$ at which consecutive classes of the second law meet, and a map
$\ell'\colon\bbN\times[0,1)\to\bbN_0$, nondecreasing in its first argument, presenting the
second law's classes against the first with a controlled ratio at every level.
\end{lemma}

\begin{proof}
\leanok
\uses{thm:chain-classes}
The identity $\theta_1=(\theta_1')^\gamma$ shows that the tail of the second geometric law at
\[
 t_k=1+\gamma(D^k-1)
\]
equals the tail of the first law at $D^k$. The numbers $t_k$ are not generally integral. We
therefore round each $t_k$ to neighbouring integers and use $U$ to divide the single atom crossed
by the cut in the required proportion. This produces consecutive integer intervals with exactly
the class masses $p^{(D)}_k$. Their endpoints remain within a fixed multiplicative factor of the
corresponding powers $D^k$, because the rounding error is at most one and
$\gamma(D-1)\geq1$. Hence the resulting map $\ell'$ has the stated monotonicity and ratio
control.
\end{proof}

\section{The bushy regime}

Drop the assumption $\theta_0=0$.

\begin{proposition}[Harris decomposition]
\label{thm:harris}
\lean{BranchingProcess.survivalMeasure_decorations_treeLaw}
\leanok
\uses{def:gw}
Let $\theta$ be supported on $\set{0,1,2}$ and supercritical with $\theta_0>0$, and let
$\cT(\omega)$ be a Galton--Watson tree with law $\theta$ conditioned on survival. Then
$\cT(\omega)$ splits as a skeleton carrying the reduced law, decorated by bushes that are
independent given the skeleton and distributed as Galton--Watson trees under the conjugate,
subcritical law. Both halves hold as equalities of laws on trees, together with their joint
form.
\end{proposition}

\begin{proof}
\leanok
The extinction probability is the root in $[0,1)$ of
$s=\theta_0+\theta_1s+\theta_2s^2$. Since the fixed-point equation factors as
$(s-1)(\theta_2s-\theta_0)=0$, this root is $q=\theta_0/\theta_2$. Conditional on the
offspring number at a vertex, its child subtrees survive independently with probability $1-q$.
Retaining the children whose subtrees survive gives the reduced skeleton. The discarded child
subtrees are independent Galton--Watson trees conditioned to become extinct, whose offspring law
is the conjugate law. Repeating this decomposition at every surviving vertex yields both the
marginal equalities of laws and their joint form.
\end{proof}

\begin{definition}[Shapes]
\label{def:shape}
\lean{ChainClasses.Shape}
\leanok
\uses{thm:harris}
A \emph{shape} is a pair $\sigma=(m,(b_1,\dots,b_{m-1}))$ with $m\geq1$ and each $b_i$ either
empty or a finite tree with offspring numbers in $\set{0,1,2}$. Its realisation is obtained
from a path $v_1\cdots v_m$ by attaching the root of each nonempty $b_i$ by an edge at $v_i$.
The entry is $v_1$, the exit is $v_m$, and $\abs\sigma$ is the number of vertices. Write $\cS$
for the countable set of shapes.
\end{definition}

\begin{proposition}[Shape decomposition]
\label{thm:shape-iid}
\lean{ChainClasses.shape_decomposition}
\leanok
\uses{def:shape}
Conditioned on infinite diameter, the shapes $S(w)$, $w\in\bbB$, are i.i.d.\ with the law
$\mu$ of $(m,(B_1,\dots,B_{m-1}))$, where $\bbP(m=k)=\widetilde\theta_1^{\,k-1}
\widetilde\theta_2$ and, given $m$, the decorations are independent, each empty with
probability $\theta_1/\widetilde\theta_1$ and otherwise a Galton--Watson tree with the
conjugate law. Moreover $\cT(\omega)$ is isometric to the assembly of the shapes along
$\bbB$.
\end{proposition}

\begin{proof}
\leanok
\uses{thm:harris, thm:geometric}
The reduced skeleton has offspring in $\set{1,2}$, so \cref{thm:geometric} makes its chain
lengths independent and geometric with ratio $\widetilde\theta_1$. By
\cref{thm:harris}, the dying bushes are conditionally independent given the skeleton. At a unary
skeleton vertex, the unused child position is empty or contains one such bush with the stated
probabilities. These choices use disjoint parts of the original offspring field for different
branch points, so the resulting shapes are i.i.d. with law $\mu$. Gluing each exit to the two
entries indexed by its children reconstructs every edge of the original tree and gives the final
isometry.
\end{proof}

\begin{lemma}[Shape mass bounds]
\label{thm:shape-mass}
\lean{ChainClasses.shape_mass_tail}
\leanok
There are $c,C,n_0>0$ depending on $\theta$ such that $\mu(\abs S\geq n)\leq e^{-cn}$ for all
$n\geq n_0$, and $\mu(\sigma)\geq e^{-C\abs\sigma}$ for every $\sigma\in\cS$ with
$\mu(\sigma)>0$.
\end{lemma}

\begin{proof}
\leanok
\uses{thm:shape-iid}
Let $p$ be the smallest positive probability among the finitely many local offspring and
decoration choices. Specifying a shape of size $n$ requires at most $2n$ such choices, so every
supported shape of size $n$ has mass at least $p^{2n}=e^{-C n}$ with
$C=2\log(p^{-1})$.

For the upper tail, the neck length has an exponential moment because it is geometric. A bush
under the conjugate law is subcritical with bounded offspring, so its total progeny also has an
exponential moment. The size of a shape is a geometric sum of these bush sizes. Choosing the
moment parameter sufficiently small gives a finite exponential moment for $\abs S$, and Markov's
inequality gives the asserted exponential tail.
\end{proof}

\section{The net quantisation}

\begin{definition}[The net and the label graph]
\label{def:shape-net}
\lean{ChainClasses.shapeFamily}
\leanok
\uses{def:shape}
Fix an integer $D\geq2$ and enumerate $\cS$ by size, breaking ties arbitrarily. Let
$\cR_D\subseteq\cS$ consist of those shapes admitting no $D$-marked quasi-isometry to any
earlier member of $\cR_D$, and for $\sigma\in\cS$ let $\rep_D(\sigma)$ be the earliest member
of $\cR_D$ to which $\sigma$ admits a $D$-marked quasi-isometry. The label graph $G_D$ has
vertex set $\cR_D$, with an edge between distinct $\sigma^*,\tau^*$ whenever they admit a
$27D^4$-marked quasi-isometry.
\end{definition}

\begin{lemma}[Basic net properties]
\label{thm:shape-net}
\lean{ChainClasses.markedQI_of_repIdx_adjacent}
\leanok
\uses{def:shape-net}
The composition of a $K$-marked and a $K'$-marked quasi-isometry is $3KK'$-marked, and every
$K$-marked quasi-isometry has a $3K^2$-marked quasi-inverse. Consequently $\rep_D(\sigma)$
exists for every $\sigma$, fixes $\cR_D$ pointwise, satisfies
$\abs{\rep_D(\sigma)}\leq\abs\sigma$, any two members of one class admit a $9D^3$-marked
quasi-isometry, and members of equal or adjacent classes admit a $729D^7$-marked
quasi-isometry.
\end{lemma}

\begin{proof}
\leanok
\uses{def:qi}
Composing the inequalities in \cref{def:qi} gives multiplicative error $KK'$, additive error
at most $2KK'$, and density and mark errors at most $KK'+2K'$. Since $K,K'\geq1$, all four
are bounded by $3KK'$. Choosing a nearest preimage of each target point similarly constructs a
$3K^2$-marked quasi-inverse.

The shape itself is always an admissible representative candidate, so $\rep_D(\sigma)$ exists.
Its definition makes it no larger than $\sigma$ and fixes representatives. To compare two shapes
with the same representative, compose the map from the first shape to the representative with a
quasi-inverse for the second. For adjacent representatives, insert the map defining the edge of
$G_D$. The preceding composition bounds give the constants $9D^3$ and $729D^7$.
\end{proof}

\begin{lemma}[Connectivity of the label graph]
\label{thm:shape-connected}
\lean{ChainClasses.netLink_connected_shapeFamily}
\leanok
\uses{def:shape-net}
The graph $G_D$ is connected.
\end{lemma}

\begin{proof}
\leanok
Let $\sigma^*\in\cR_D$ have more than one vertex. Delete a leaf of one of its bushes if possible,
and otherwise delete the exit of its neck. The remaining tree is the realisation of a shape
$\tau$ with $\abs\tau=\abs{\sigma^*}-1$. The map that fixes the remaining vertices is a
$2$-marked quasi-isometry from $\sigma^*$ to $\tau$. Composing with the defining map from
$\tau$ to $\rep_D(\tau)$ shows that $\sigma^*$ is adjacent to the strictly smaller
representative $\rep_D(\tau)$. Induction on the size connects every representative to the
one-vertex shape.
\end{proof}

\begin{lemma}[Entropy dilution]
\label{thm:dilution}
\lean{ChainClasses.dilution}
\leanok
\uses{def:shape-net}
There is an absolute constant $c_0$ such that for all integers $D\geq2$ and $n\geq1$, every
family of shapes of size at most $n$, no two of which admit a $D$-marked quasi-isometry, has
at most $\exp\bigl(c_0(nD^{-1/3}+1)\bigr)$ members. In particular there are at most that many
representatives of size at most $n$.
\end{lemma}

\begin{proof}
\leanok
Encode a finite rooted tree by its depth-first traversal, writing $1$ on entering a vertex and
$2$ on leaving it. A tree with $k$ vertices has a code of length $2k$, and the running excess
of $1$'s over $2$'s shows that the code determines the rooted tree. Adding the marked exit costs
at most a factor $k$, so the number of rooted trees with a marked vertex and at most $k$ vertices
is at most $e^{3k}$.

Now cut a shape into connected parts of size comparable to a scale $s$. Each complete part has
at least $s$ vertices and every part has fewer than $2s$ vertices. Contracting the parts gives a
rooted tree with a marked vertex and at most $n/s+1$ vertices, and the contraction map is
$2s$-marked. If two shapes have isomorphic contractions, composing one contraction with a
quasi-inverse of the other gives a $72s^3$-marked quasi-isometry between them. Choose
$s\asymp D^{1/3}$ so that $72s^3\leq D$. Pairwise $D$-incomparability then forces the contracted
trees to be distinct. The depth-first count is consequently at most
$\exp\bigl(c_0(nD^{-1/3}+1)\bigr)$. For the bounded range of smaller $D$, the same estimate
follows by taking $s=1$ and increasing the absolute constant.
\end{proof}

\section{The glued transfer}

\begin{lemma}[Glued transfer]
\label{thm:glued-transfer}
\lean{ChainClasses.Assembly.glued_transfer}
\leanok
\uses{def:qi}
Let $(\sigma_w)_{w\in\bbB}$ and $(\sigma'_w)_{w\in\bbB}$ be families of shapes with assemblies
$\cT$ and $\cT'$, let $K\geq1$, and suppose that for every $w\in\bbB$ the shapes $\sigma_w$
and $\sigma'_w$ admit a $K$-marked quasi-isometry. Then there is an $8K^2$-quasi-isometry
$\cT\to\cT'$.
\end{lemma}

\begin{proof}
\leanok
\uses{thm:shape-iid}
Choose a $K$-marked quasi-isometry between each pair $\sigma_w$ and $\sigma'_w$, and define
$\psi$ copy by copy. Each copy is an induced subtree of its assembly, so its internal distances
are unchanged. The marked conditions ensure that the entry and exit of a source copy map within
distance $K$ of the corresponding marks in the target copy.

A geodesic between different copies consists of at most two partial paths inside its endpoint
copies and a sum of complete necks between them. The marked quasi-isometries compare each partial
path and compare every complete neck in both directions. Summing these inequalities along the
geodesic gives distortion at most $8K^2$ in the ancestor and divergence cases. Coarse density
holds copy by copy, since every point of $\sigma'_w$ lies within distance $K$ of the image of the
map on $\sigma_w$. Hence $\psi$ is an $8K^2$-quasi-isometry.
\end{proof}
